\usepackage[english,italian]{babel}
\usepackage{amsmath, inconsolata, dsfont, mathrsfs, amsthm, amssymb, pifont, tikz, hyperref, newlfont, mathabx, mathtools, graphicx, caption, array, multirow, tabularx, svg, pdfpages, setspace}
\usepackage{listings, float, aliascnt, caption}
\usepackage{tikz}
\usepackage[T1]{fontenc}
\usepackage[utf8]{inputenc}
\usepackage[]{geometry}
\renewcommand{\lstlistingname}{Listato}
\usepackage{xcolor}
\usepackage{setspace}
\usepackage{csquotes}
\usepackage[utf8]{inputenc}
\usepackage{listings}
%\usepackage[
%backend=bibtex,
%style=numeric,
%sorting=ynt
%]{biblatex}
\usetikzlibrary{arrows.meta}

\hypersetup{
bookmarksopen=true,	
pdfauthor=Giacomo Borra
}

\lstset{
    literate=
        {à}{{\`a}}1
        {è}{{\`e}}1
        {é}{{\'e}}1
        {ì}{{\`i}}1
        {ò}{{\`o}}1
        {ù}{{\`u}}1
}

\lstdefinestyle{myStile}{
    language=Python,
    basicstyle=\ttfamily\scriptsize,
    numbers=left,
    numberstyle=\tiny\color{gray},
    numbersep=5pt,
    frame=single,
    breaklines=true,
    showstringspaces=false,
    aboveskip=0.8em,
    belowskip=0.8em
}

\theoremstyle{definition}

\newtheorem{proposition}{Proposizione}[section]

\newaliascnt{theorem}{proposition}
\newtheorem{theorem}[theorem]{Teorema}
\aliascntresetthe{theorem}

\newaliascnt{lemma}{proposition}
\newtheorem{lemma}[lemma]{Lemma}
\aliascntresetthe{lemma}

\newaliascnt{definition}{proposition}
\newtheorem{definition}[definition]{Definizione}
\aliascntresetthe{definition}

\newaliascnt{corollary}{proposition}
\newtheorem{corollary}[corollary]{Corollario}
\aliascntresetthe{corollary}

\newaliascnt{example}{proposition}
\newtheorem{example}[example]{Esempio}
\aliascntresetthe{example}

\newaliascnt{remark}{proposition}
\newtheorem{remark}[remark]{Remark}
\aliascntresetthe{remark}

\numberwithin{equation}{section}

\providecommand*{\propositionautorefname}{Proposizione}
\providecommand*{\theoremautorefname}{Teorema}
\providecommand*{\lemmaautorefname}{Lemma}
\providecommand*{\definitionautorefname}{Definizione}
\providecommand*{\corollaryautorefname}{Corollario}
\providecommand*{\exampleautorefname}{Esempio}
\providecommand*{\remarkautorefname}{Remark}
\providecommand*{\equationautorefname}{Equazione}

\newcommand{\norm}[1]{\left\|#1\right\|}
\renewcommand{\sectionautorefname}{§}
\renewcommand{\subsectionautorefname}{§}
\renewcommand{\subsubsectionautorefname}{§}